\section{Avaluació}
\label{cap:evaluation}

Aquest capítol presenta l’avaluació del prototip a partir d’una bateria d’escenaris controlats i reproduïbles. L’objectiu és comprovar si el sistema compleix els requisits definits al Capítol~\ref{cap:design} i si és capaç d’oferir assistència útil abans de la confirmació d’una transacció dins de MetaMask.

L’avaluació se centra en quatre aspectes principals: la intercepció de transaccions, la detecció de patrons de risc, la generació d’explicacions comprensibles i la integració completa entre la DApp, el Snap, el backend, la base de dades local, Ollama i els contractes desplegats a Sepolia.

Cada escenari s’exposa conjuntament amb el resultat esperat, el resultat observat i els requisits que permet validar. Aquesta organització permet relacionar de manera directa cada prova amb el comportament del sistema i amb el grau de compliment dels requisits.

\subsection{Infraestructura i preparació de l’avaluació}

Abans de presentar els escenaris de prova i els resultats obtinguts, cal descriure la infraestructura utilitzada per executar l’avaluació i el procediment seguit per preparar les operacions analitzades. Aquesta preparació permet garantir que les proves es duen a terme en un entorn controlat, reproduïble i prou proper a una interacció real amb la xarxa Ethereum. 

L’entorn d’avaluació combina els components locals del prototip —la DApp de proves, MetaMask Flask, el Snap, el backend d’anàlisi, la base de dades SQLite i el servei local d’Ollama— amb una infraestructura blockchain externa basada en la xarxa de proves Sepolia. 

Sobre aquesta xarxa es despleguen contractes intel·ligents específicament dissenyats per generar operacions representatives, com aprovacions ERC20 limitades o il·limitades i permisos globals sobre actius ERC721. 

Aquesta secció descriu el disseny dels contractes de prova, la connexió amb Sepolia, el procés de compilació i desplegament, la verificació mitjançant Etherscan i el protocol utilitzat per generar les transaccions i recollir-ne les evidències. L’objectiu és establir les condicions experimentals necessàries abans d’analitzar, en els apartats posteriors, el comportament del sistema en cada escenari.

\subfile{impl_contracts}

\subfile{protocol_execucio}

\subfile{monitoritzacio_avaluacio}

\subfile{tracabilitat_requisits}

\subfile{escala_risc}

\subfile{escenaris_resultats}

\subsection{Anàlisi dels resultats}
\label{sec:analisi_resultats}

Un cop presentats els escenaris i les evidències recollides, aquest apartat analitza els resultats des de tres perspectives complementàries. En primer lloc, es quantifica la capacitat de classificació del sistema mitjançant una matriu de confusió i les mètriques derivades. A continuació, es contrasta cada requisit definit al Capítol~\ref{cap:design} amb les evidències obtingudes. Finalment, es delimita l’abast dels resultats i s’interpreten els casos límit detectats.

\subfile{matriu_confussio}

\subfile{compliment_requisits}

\subfile{limitacions_avaluacio}
